\section{共享年表之二：改革竞赛变成吞并竞赛}

\begin{event}{公元前334年}{齐魏互称王}{可靠纪年}{徐州一带}\eventanchor{WQ-0334-KINGS}
\term{这一年发生了什么}：齐、魏相互承认“王”号。
\term{后来改变了什么}：周天子的名分继续被各国使用，却不再能限制最高称号的竞争。战国国家公开把自身放到接近“王”的位置上。
\end{event}

\begin{event}{公元前307年}{赵武灵王胡服骑射}{可靠纪年}{赵}\eventanchor{WQ-0307-HUFU}
\term{这一年发生了什么}：赵武灵王推动服饰和骑射改革。
\term{事情为什么重要}：这不是单纯换衣服，而是一个中原国家为了面对北方边地战争和机动骑兵，愿意调整旧有身份象征与军事技术。
\end{event}

\begin{event}{公元前293年}{伊阙之战}{可靠纪年}{伊阙}\eventanchor{WQ-0293-YIQUE}
\term{这一年发生了什么}：秦军在伊阙重创韩魏联军。
\term{后来改变了什么}：关中秦国向东推进的通道被进一步打开，韩魏的缓冲作用下降。
\end{event}

\begin{event}{公元前260年}{长平之战}{可靠纪年}{长平}\eventanchor{WQ-0260-CHANGPING}
\term{这一年发生了什么}：秦赵长期对峙后，赵军败降；关于降卒处置的规模，传世叙事极具震撼力，也应注意其数字与修辞的审视。
\term{后来改变了什么}：赵国精锐受重创，东方抗秦均衡被打破。秦的统一不再只是可能性，而开始进入可操作的军事步骤。
\end{event}

\begin{event}{公元前256年}{东周亡}{可靠纪年}{洛邑}\eventanchor{WQ-0256-ZHOUDIE}
\term{这一年发生了什么}：秦灭东周残余政权。
\term{后来改变了什么}：周天子作为名义共主的最后政治实体消失。六国仍在，但“天下”已经没有一个可供各方共同借用的旧中心。
\end{event}
